\paragraph{Overall and reference-domain comparisons.}
Across the 15 confirmation conditions, the full ensemble improves on the
ungated residual in 14 and on the binary gate in 13. Its advantage is largest
on the two-dimensional FEM reference: it improves all five conditions against
both controls, reducing NRMSE by 18.1--73.3\% relative to the binary gate.
For example, in-support random NRMSE is 0.0103 for the full model, 0.0387 for
the binary gate and 0.0529 for the ungated residual. These differences support
the utility of the full hybrid design under the reference-domain training
recipe, rather than attributing all gains to cross-table learning alone.

Transfer to ASAP7 is less uniform. Full routing improves on the ungated
residual in all five conditions, but the binary ensemble is better on
out-of-support block masks (0.1291 versus 0.1358) and extrapolation masks
(0.1037 versus 0.1073). The full branch bank is consequently not a uniformly
better choice across the evaluated settings.

\paragraph{Public-rule attribution.}
The full ensemble has lower point estimates than the binary gate in all five
public conditions, but the relative improvements range from only 0.6\% to
5.6\%. The ungated residual has a lower point estimate for in-support block masks
(0.0645 versus 0.0678), while the full model is better in the other four
conditions. Much of the reduction relative to numerical interpolation is
therefore already achieved by a support-conditioned residual Transformer.

Table~\ref{tab:control-bootstrap} qualifies the small routing gains. The full
model's paired intervals exclude zero against the binary gate in the two
in-support conditions, but not in any of the three out-of-support conditions.
Against the ungated residual, they exclude zero in favor of the full model
for in-support random and out-of-support block masks. These are pointwise
comparisons conditional on the trained ensembles, not evidence of a uniform
advantage across training seeds or a guarantee that the full branch bank is
necessary. Together with the pruning diagnostic, the public evidence supports
a strong learned residual and more limited gains from the extra routing machinery.

\paragraph{Training-seed checks and interpretation.}
Replaying all original full-model members reproduces the reported ensemble
NRMSEs exactly in all 15 conditions. Comparing corresponding individual seeds,
the full model beats the ungated residual in 12/15 public, 15/15 reference and
13/15 transfer seed--condition pairs; against the binary gate the counts are
9/15, 13/15 and 13/15. These descriptive counts are not independent statistical
tests. They also need not match ensemble win counts, because averaging
predictions precedes computing ensemble NRMSE. The reference advantage is
stronger than the public advantage, but neither seed counts nor pooled
ensemble scores establish uniform dominance. In particular, these controls
support a condition-dependent benefit of the full hybrid design, not the
necessity of every analytic operator in every domain.
